% Standalone appendix fragment for the retaliation-capacity cascade.
% Compile by wrapping this file in a LaTeX document that defines theorem
% environments and calls \appendix before \input.

\section{Discrete Sobolev Compactness and Weak Convergence}
\label{app:discrete-cascade-compactness}

This appendix isolates the analytic step behind the retaliation-capacity
cascade.  The discrete cascade is treated as a sequence of Dirichlet
Sturm--Liouville problems.  As established by the bridge argument in the main text, the near-null discrete fields are not an ad hoc assumption but the direct consequence of the game's finite cascade reaching a dynamic quasi-equilibrium. The point is not to impose an Airy profile by
analogy, but first to prove that these near-null discrete fields have a continuum
weak limit for the Sturm--Liouville operator.

Let $I=[a,b]$ and let $h=(b-a)/(n+1)$ with grid points
$x_i=a+ih$, $0\le i\le n+1$.  Define the discrete Dirichlet space
\[
  X_h
  :=
  \{u=(u_0,\ldots,u_{n+1}) : u_0=u_{n+1}=0\}.
\]
For $u,v\in X_h$, set
\[
  \langle u,v\rangle_{L_h^2}
  :=
  h\sum_{i=1}^n u_i v_i,
  \qquad
  \|u\|_{L_h^2}^2
  :=
  \langle u,u\rangle_{L_h^2},
\]
and define the backward difference
\[
  (D_hu)_i := \frac{u_i-u_{i-1}}{h},
  \qquad 1\le i\le n+1.
\]
The discrete $H^1$ norm is
\[
  \|u\|_{H_h^1}^2
  :=
  \|u\|_{L_h^2}^2
  +
  h\sum_{i=1}^{n+1}|(D_hu)_i|^2 .
\]
The discrete Laplacian is
\[
  (\Delta_hu)_i
  :=
  \frac{u_{i+1}-2u_i+u_{i-1}}{h^2},
  \qquad 1\le i\le n,
\]
with the boundary values $u_0=u_{n+1}=0$.  Given $\sigma>0$ and a bounded
grid potential $V_h=(V_{h,i})_{i=1}^n$, define
\[
  (H_hu)_i := -\sigma(\Delta_hu)_i+V_{h,i}u_i.
\]
The dual residual is measured by
\[
  \|r_h\|_{H_h^{-1}}
  :=
  \sup_{\varphi\in X_h,\ \|\varphi\|_{H_h^1}=1}
  \left|
    h\sum_{i=1}^n r_{h,i}\varphi_i
  \right|.
\]

Let $I_hu$ be the continuous piecewise-linear interpolant of $u$ with nodal
values $u_i$ at $x_i$.  The elementary interpolation estimates used below are
\[
  c_0\|u\|_{H_h^1}
  \le
  \|I_hu\|_{H_0^1(I)}
  \le
  C_0\|u\|_{H_h^1},
  \tag{A.1}
\]
with constants independent of $h$.  The Dirichlet boundary also gives a
uniform discrete Poincare inequality,
\[
  \|u\|_{L_h^2}
  \le
  C_P
  \left(h\sum_{i=1}^{n+1}|(D_hu)_i|^2\right)^{1/2}.
  \tag{A.2}
\]
If one works with Neumann data or conditional measures, the analogous statement
is a quotient estimate modulo constants; here the zero boundary data remove
that quotient by fixing the centered representative.

\begin{theorem}[Discrete Cascade-to-Airy Convergence]
\label{thm:discrete-cascade-to-airy}
Let $\sigma>0$.  Suppose $V\in C([a,b])$ and the grid potentials satisfy
\[
  \max_{1\le i\le n}|V_{h,i}-V(x_i)|\to 0,
  \qquad
  \sup_{h,i}|V_{h,i}|<\infty .
\]
Let $\psi_h\in X_h$ be a sequence of discrete retaliation-capacity fields such
that
\[
  \|\psi_h\|_{L_h^2}=1,
  \qquad
  \|H_h\psi_h\|_{H_h^{-1}}\to 0 .
\]
Then, after passing to a subsequence, $I_h\psi_h$ converges weakly in
$H_0^1(I)$ and strongly in $L^2(I)$ to a non-zero field
$\Psi\in H_0^1(I)$ with $\|\Psi\|_{L^2(I)}=1$.  The limit solves the weak
Sturm--Liouville equation
\[
  \sigma\int_a^b \Psi'(x)\eta'(x)\,dx
  +
  \int_a^b V(x)\Psi(x)\eta(x)\,dx
  =
  0
  \qquad
  \text{for every }\eta\in H_0^1(I).
  \tag{A.3}
\]
Equivalently, $-\sigma\Psi''+V\Psi=0$ in $H^{-1}(I)$.  If, in addition,
$V$ is $C^2$ near a simple turning point, this weak limit has the Airy normal
form described in the \hyperref[app:airy-normal-form]{Airy normal-form appendix}.
\end{theorem}

\begin{proof}
Write the discrete bilinear form associated with $H_h$ as
\[
  a_h(u,\varphi)
  :=
  \sigma h\sum_{i=1}^{n+1}(D_hu)_i(D_h\varphi)_i
  +
  h\sum_{i=1}^n V_{h,i}u_i\varphi_i .
\]
Discrete summation by parts gives
\[
  a_h(u,\varphi)
  =
  h\sum_{i=1}^n (H_hu)_i\varphi_i,
  \qquad u,\varphi\in X_h .
  \tag{A.4}
\]
The potential may change sign, so the unshifted form is not assumed coercive.
Choose $\gamma>1+\sup_{h,i}|V_{h,i}|$ and define the shifted form
\[
  a_h^\gamma(u,u)
  :=
  a_h(u,u)+\gamma\|u\|_{L_h^2}^2 .
\]
Then, with $c_\gamma:=\min\{\sigma,\gamma-\sup_{h,i}|V_{h,i}|\}>0$,
\[
  a_h^\gamma(u,u)
  \ge
  c_\gamma h\sum_{i=1}^{n+1}|(D_hu)_i|^2
  +
  c_\gamma\|u\|_{L_h^2}^2 ,
  \tag{A.5}
\]
This is the Garding shift: the positive shift supplies coercivity without
claiming that the sign-changing operator $H_h$ is coercive on its own.

Using (A.4), the residual bound, and $\|\psi_h\|_{L_h^2}=1$,
\[
  a_h^\gamma(\psi_h,\psi_h)
  =
  h\sum_{i=1}^n (H_h\psi_h)_i\psi_{h,i}
  +
  \gamma
  \le
  \|H_h\psi_h\|_{H_h^{-1}}\|\psi_h\|_{H_h^1}
  +
  \gamma .
\]
Together with (A.5), this quadratic inequality gives a uniform
$H_h^1$ bound on $\psi_h$.  By the interpolation equivalence (A.1), the
piecewise-linear fields $I_h\psi_h$ are uniformly bounded in $H_0^1(I)$.
Rellich compactness therefore gives a subsequence, not relabeled, and
$\Psi\in H_0^1(I)$ such that
\[
  I_h\psi_h\rightharpoonup \Psi \quad\text{in }H_0^1(I),
  \qquad
  I_h\psi_h\to \Psi \quad\text{in }L^2(I).
\]
The normalization transfers to the limit because the discrete $L_h^2$ norm and
the $L^2$ norm of the piecewise-linear interpolant differ by $o(1)$ on
uniformly $H_h^1$-bounded sequences.  Hence $\|\Psi\|_{L^2(I)}=1$.

It remains to pass to the weak equation.  Fix $\eta\in C_c^\infty(a,b)$ and
let $\eta_h\in X_h$ be its nodal restriction, $\eta_{h,i}=\eta(x_i)$.  Since
$\|\eta_h\|_{H_h^1}$ is uniformly bounded, the residual assumption gives
\[
  a_h(\psi_h,\eta_h)
  =
  h\sum_{i=1}^n (H_h\psi_h)_i\eta(x_i)
  \to 0 .
  \tag{A.6}
\]
For the gradient part,
\[
  h\sum_{i=1}^{n+1}(D_h\psi_h)_i(D_h\eta_h)_i
  =
  \int_a^b (I_h\psi_h)'(x)(I_h\eta_h)'(x)\,dx
  \to
  \int_a^b \Psi'(x)\eta'(x)\,dx,
\]
using weak convergence in $H^1$ and the strong convergence
$I_h\eta_h\to\eta$ in $H_0^1(I)$.  For the potential term, uniform convergence
of $V_h$ to $V$, strong $L^2$ convergence of $I_h\psi_h$, and the consistency
of nodal quadrature for the smooth test function give
\[
  h\sum_{i=1}^n V_{h,i}\psi_{h,i}\eta(x_i)
  \to
  \int_a^b V(x)\Psi(x)\eta(x)\,dx .
\]
Taking limits in (A.6) proves (A.3) for smooth compactly supported tests.
Density of $C_c^\infty(a,b)$ in $H_0^1(I)$ extends the identity to every
$\eta\in H_0^1(I)$.
\end{proof}

\section{Airy Normal Form near a Simple Turning Point}
\label{app:airy-normal-form}

The Airy profile appears only after the continuum weak equation has been
obtained.  Assume that $V\in C^2$ in a neighborhood of $x_c$ and that
\[
  V(x_c)=0,
  \qquad
  V'(x_c)=-c<0 .
  \tag{B.1}
\]
By one-dimensional elliptic regularity applied to
$-\sigma\Psi''+V\Psi=0$, the weak limit from
Theorem~\ref{thm:discrete-cascade-to-airy} is $C^2$ locally near $x_c$.
Taylor expansion gives
\[
  V(x)=-c(x-x_c)+O((x-x_c)^2).
\]
Set
\[
  z := \left(\frac{c}{\sigma}\right)^{1/3}(x-x_c),
  \qquad
  Y(z):=\Psi\!\left(x_c+\left(\frac{\sigma}{c}\right)^{1/3}z\right).
\]
Then the local equation becomes
\[
  Y''(z)+zY(z)=O(z^2)Y(z).
  \tag{B.2}
\]
Keeping the first-order turning-point term gives the Airy normal form
\[
  Y''(z)+zY(z)=0.
  \tag{B.3}
\]
Its local solution space is two-dimensional:
\[
  Y(z)=A\,\operatorname{Ai}(-z)+B\,\operatorname{Bi}(-z).
  \tag{B.4}
\]

Equation (B.4) is a local normal form, not a branch selection rule.  The pure
$\operatorname{Ai}(-z)$ branch is justified only after an additional matching
condition is imposed, for example finite-energy matching from the protected
side or a boundary condition that eliminates the growing
$\operatorname{Bi}(-z)$ component.  Without that matching information, the
discrete cascade theorem identifies the continuum Sturm--Liouville limit and
the local Airy basis, but it does not select $B=0$.
